\definecolor{GreenD}{HTML}{3B6E10}
\definecolor{GreenL}{HTML}{E3F2D5}

\begin{tikzpicture}[
    font=\small,
    >={Stealth[length=2.2mm, width=1.5mm]},
    %--- tensor de entrada (cinza, igual a arquitetura.tex)
    tensor/.style={
        draw=gray!55, thick, rounded corners=8pt,
        fill=gray!9, minimum width=2.8cm, minimum height=0.65cm,
        align=center},
    %--- bloco de rede (verde chamfered, igual a arquitetura.tex)
    net/.style={
        draw=GreenD, very thick, chamfered rectangle,
        chamfered rectangle xsep=0.22cm, fill=GreenL,
        minimum width=2.6cm, minimum height=0.65cm,
        align=center},
    %--- variável intermediária (verde tracejado)
    rv/.style={
        draw=GreenD, thick, dashed, rounded corners=8pt,
        fill=GreenL!50, minimum width=2.6cm, minimum height=0.65cm,
        align=center},
]

%% ── ENTRADAS (fora da caixa) ─────────────────────────────────
% e entra pelo topo; z entra pela esquerda no nível de ⊗
\node[tensor] (e)       at ( 0.0, 7.6) {Embeddings $e$};
\node[tensor] (factors) at (-5.0, 3.8) {Fatores $z$};

%% ── RAMO BETA (esquerda) E RAMO ALPHA (direita) ─────────────
\node[net] (betalayer)  at (-1.9, 6.25) {Beta Layer $\varphi_{\beta}$};
\node[net] (alphalayer) at ( 1.9, 6.25) {Alpha Layer $\pi_{\alpha}$};

\node[rv] (beta)  at (-1.9, 5.05) {Exposições $\beta$};
\node[rv] (alpha) at ( 1.9, 5.05) {Componente $\alpha$};

%% ── OPERADORES ───────────────────────────────────────────────
\node[draw=GreenD, very thick, circle,
      minimum size=0.78cm, fill=white, inner sep=0pt, font=\normalsize]
    (otimes) at (-1.9, 3.8) {$\otimes$};

\node[draw=GreenD, very thick, circle,
      minimum size=0.78cm, fill=white, inner sep=0pt, font=\normalsize]
    (oplus) at (0.0, 2.6) {$\oplus$};

%% ── SAÍDA (fora da caixa, base) ─────────────────────────────
\node[rv, fill=GreenL, draw=GreenD, very thick] (yhat) at (0, 0.35)
    {Retorno reconstruído $\hat{y}$};

%% ── CAIXA DO MÓDULO ─────────────────────────────────────────
\draw[GreenD, thick, rounded corners=12pt]
    (-3.45, 0.85) rectangle (3.45, 6.85);

%% ── SETAS ───────────────────────────────────────────────────
% e bifurca para os dois ramos
\coordinate (branch) at ($(e.south) + (0, -0.3)$);
\draw[gray!60, thick]         (e.south) -- (branch);
\draw[->, gray!60, thick]     (branch)  -| (betalayer.north);
\draw[->, gray!60, thick]     (branch)  -| (alphalayer.north);

% z entra horizontalmente no ⊗
\draw[->, gray!60, thick]     (factors.east)  -- (otimes.west);

% Fluxo interno (verde)
\draw[->, GreenD, thick]      (betalayer.south)  -- (beta.north);
\draw[->, GreenD, thick]      (alphalayer.south) -- (alpha.north);
\draw[->, GreenD, thick]      (beta.south)       -- (otimes.north);
\draw[->, GreenD, thick]      (otimes.south)     |- (oplus.west);
\draw[->, GreenD, thick]      (alpha.south)      |- (oplus.east);

% Saída
\draw[->, GreenD, thick]      (oplus.south)      -- (yhat.north);

%% ── RÓTULO ──────────────────────────────────────────────────
\node[font=\small, anchor=north] at (0, -0.3)
    {(b) Factor Decoder};

\end{tikzpicture}